% Shared preamble of the three slide decks of this unit
% (slides_lm1.tex, slides_lm2.tex, slides_lm3.tex).
% The look is the same as the PINN and CNN slides.
% The figure files in tikz/ are the ones of the notes (main.tex),
% so the colors and TikZ styles below are the same as in the notes.

% -------------------------------------------------
% Packages
% -------------------------------------------------
\usepackage{kotex}
\usepackage{fontspec}
\usepackage{amsmath,amssymb,bm,mathtools}
\usepackage{booktabs,array}
\usepackage{xcolor}
\usepackage{colortbl}
\usepackage{tikz}
\usetikzlibrary{positioning,arrows.meta,calc,fit,backgrounds,
  decorations.pathreplacing,shapes.misc}
\usepackage[most]{tcolorbox}
\usepackage{listings}
\usepackage{adjustbox}

% -------------------------------------------------
% Colors (same as the notes)
% -------------------------------------------------
\definecolor{deepblue}{RGB}{34,71,120}
\definecolor{softblue}{RGB}{239,246,253}
\definecolor{softgray}{RGB}{246,247,249}
\definecolor{softgreen}{RGB}{239,249,242}
\definecolor{softyellow}{RGB}{255,249,229}
\definecolor{codegray}{RGB}{248,248,248}
% colors of the figures
\definecolor{tokc}{RGB}{31,100,180}
\definecolor{hotc}{RGB}{200,30,40}
\definecolor{attc}{RGB}{214,110,0}
\definecolor{grnc}{RGB}{40,140,70}

% TikZ styles used by the files in tikz/ (same as the notes)
\tikzset{
  >=Latex,
  tok/.style={draw=#1!70!black,fill=#1!10,rounded corners=2pt,font=\ttfamily,
    minimum height=7mm,inner sep=3pt},
  tok/.default=tokc,
  lay/.style={draw=#1!70!black,fill=#1!8,rounded corners=3pt,align=center,
    minimum height=8mm,minimum width=22mm,inner sep=3pt},
  lay/.default=black,
  neu/.style={circle,draw=black!70,minimum size=5.5mm,inner sep=0pt,font=\small},
  ar/.style={->,thick,#1},
  ar/.default=black!80,
  op/.style={circle,draw,inner sep=1pt,minimum size=6mm,font=\small},
  lab/.style={font=\small,align=center},
}

% -------------------------------------------------
% Beamer look (same as the PINN and CNN slides)
% -------------------------------------------------
\usetheme{default}
\usefonttheme[onlymath]{serif}
\setbeamertemplate{navigation symbols}{}
\setbeamersize{text margin left=8mm,text margin right=8mm}
\setbeamercolor{structure}{fg=deepblue}
\setbeamercolor{frametitle}{fg=deepblue}
\setbeamercolor{title}{fg=deepblue}
\setbeamerfont{frametitle}{size=\large,series=\bfseries}
\setbeamertemplate{itemize items}[circle]
\setbeamertemplate{enumerate items}[default]
\setbeamertemplate{frametitle}{%
  \vskip3mm
  {\usebeamerfont{frametitle}\usebeamercolor[fg]{frametitle}\insertframetitle\par}%
  \vskip1mm
  {\color{deepblue!45}\hrule height 0.6pt}%
  \vskip1mm}
\setbeamertemplate{footline}{%
  \hbox to\paperwidth{%
    \hspace*{8mm}{\color{black!45}\scriptsize \insertshorttitle
      \ifx\insertsectionhead\empty\else\ \ $\cdot$\ \ \insertsectionhead\fi}%
    \hfill{\color{black!45}\scriptsize\insertframenumber\,/\,\inserttotalframenumber}%
    \hspace*{8mm}}%
  \vskip2.5mm}
\setbeamertemplate{title page}{%
  \vfill
  \begin{center}
  {\color{deepblue}\Huge\bfseries\inserttitle\par}
  \vskip3mm
  {\color{deepblue!60}\rule{9cm}{0.8pt}\par}
  \vskip3mm
  {\large\insertsubtitle\par}
  \vskip8mm
  {\normalsize\insertdate\par}
  \end{center}
  \vfill}
\AtBeginSection[]{%
  \begin{frame}
  \vfill
  \centering
  {\color{deepblue!55}\Large\bfseries \insertsectionnumber\par}
  \vskip2mm
  {\color{deepblue}\LARGE\bfseries\insertsectionhead\par}
  \vskip3mm
  {\color{deepblue!50}\rule{7cm}{0.8pt}\par}
  \vfill
  \end{frame}}
\setlength{\leftmargini}{1.2em}
\setlength{\leftmarginii}{1.1em}
% numbered sections in the table of contents. The second and third decks start at
% section 5 and 9 like the notes: \setcounter{section}{4} before the first \section,
% and \tableofcontents[firstsection=-3] (third deck: -7) so that the list shows 5-8 (9-11).
\setbeamertemplate{section in toc}[sections numbered]

% -------------------------------------------------
% Boxes (same colors as the notes). The title is braced, so it may contain commas.
% -------------------------------------------------
\newtcolorbox{keybox}[1][]{colback=softblue,colframe=deepblue,boxrule=0.7pt,arc=2mm,
  title={#1},fonttitle=\bfseries\small,left=2mm,right=2mm,top=1mm,bottom=1mm,
  toptitle=0.5mm,bottomtitle=0.5mm}
\newtcolorbox{thinkbox}[1][]{colback=softyellow,colframe=orange!70!black,boxrule=0.6pt,
  arc=2mm,title={#1},fonttitle=\bfseries\small,left=2mm,right=2mm,top=1mm,bottom=1mm,
  toptitle=0.5mm,bottomtitle=0.5mm}
\newtcolorbox{remarkbox}[1][]{colback=softgray,colframe=black!45,boxrule=0.5pt,arc=2mm,
  title={#1},fonttitle=\bfseries\small,left=2mm,right=2mm,top=1mm,bottom=1mm,
  toptitle=0.5mm,bottomtitle=0.5mm}
\newtcolorbox{mlbox}[1][]{colback=softgreen,colframe=green!45!black,boxrule=0.6pt,
  arc=2mm,title={#1},fonttitle=\bfseries\small,left=2mm,right=2mm,top=1mm,bottom=1mm,
  toptitle=0.5mm,bottomtitle=0.5mm}

% -------------------------------------------------
% Code style
% -------------------------------------------------
\lstdefinestyle{pythonstyle}{
  language=Python,
  backgroundcolor=\color{codegray},
  basicstyle=\ttfamily\scriptsize,
  keywordstyle=\bfseries\color{deepblue},
  commentstyle=\itshape\color{black!55},
  stringstyle=\color{black!75},
  numbers=none,
  frame=single,
  rulecolor=\color{black!25},
  breaklines=true,
  showstringspaces=false,
  tabsize=4,
  columns=fullflexible,
  keepspaces=true,
  aboveskip=1mm,
  belowskip=1mm
}
\lstset{style=pythonstyle}

% -------------------------------------------------
% Math commands (same as the notes)
% -------------------------------------------------
\newcommand{\R}{\mathbb{R}}
\newcommand{\vect}[1]{\bm{#1}}
\newcommand{\T}{^{\mathsf T}}
\newcommand{\loss}{\mathcal{L}}
\newcommand{\softmax}{\operatorname{softmax}}
\newcommand{\w}[1]{\texttt{#1}}

% -------------------------------------------------
% Figures
% -------------------------------------------------
% one-line caption under a figure
\newcommand{\figcap}[1]{\par\vskip1mm{\centering\footnotesize\color{black!70}#1\par}}
% TikZ figure of the notes, shrunk to fit: \sfig[max height]{file name in tikz/}
\newcommand{\sfig}[2][5.4cm]{%
  \par{\centering
  \begin{adjustbox}{max width=\linewidth,max totalheight=#1}\input{tikz/#2}\end{adjustbox}
  \par}}
% picture from figures/, fit into the given height: \pfig[max height]{file}
\newcommand{\pfig}[2][5.4cm]{%
  \par{\centering
  \includegraphics[width=\linewidth,height=#1,keepaspectratio]{#2}\par}}
